%!TEX root = ../thesis.tex
% Cтворюємо перелік умовних позначень, скорочень і термінів у форматі поняття --- розшифровка
% Скорочення потрібно подавати в алфавітному порядку. Скорочення латиницею повинні передувати скороченням кирилицею

%!TEX root = ../thesis.tex
% Cтворюємо перелік умовних позначень, скорочень і термінів у форматі поняття --- розшифровка
% Скорочення потрібно подавати в алфавітному порядку. Скорочення латиницею повинні передувати скороченням кирилицею

ADSR --- Attack, Decay, Sustain, Release --- часова огинаюча звуку: атака, затримка, утримання, спад

BCELoss --- Binary Cross-Entropy Loss --- бінарна крос-ентропійна функція втрат

BiLSTM --- Bidirectional Long Short-Term Memory --- двонапрямлена рекурентна нейронна мережа з довгою короткочасною пам'яттю

CNN --- Convolutional Neural Network --- згорткова нейронна мережа

DCT --- Discrete Cosine Transform --- дискретне косинусне перетворення

GPU --- Graphics Processing Unit --- графічний процесор, використовується для прискорення навчання нейронних мереж

IRMAS --- Instrument Recognition in Musical Audio Signals --- набір даних для розпізнавання інструментів у музичних аудіосигналах

LSTM --- Long Short-Term Memory --- рекурентна нейронна мережа з довгою короткочасною пам'яттю

mAP --- mean Average Precision --- середня точність, стандартна метрика для multi-label класифікації

MFCC --- Mel-Frequency Cepstral Coefficients --- мел-частотні кепстральні коефіцієнти

MIDI --- Musical Instrument Digital Interface --- стандартний протокол цифрового представлення музичних подій

MIREX --- Music Information Retrieval Evaluation eXchange --- міжнародний конкурс оцінювання систем музичної інформатики

ML --- Machine Learning --- машинне навчання

STFT --- Short-Time Fourier Transform --- перетворення Фур'є короткого вікна

SVM --- Support Vector Machine --- метод опорних векторів

МН --- машинне навчання

ОНД --- основи наукових досліджень
